\documentclass[11pt]{article}

% ---------- Packages ----------
\usepackage[margin=1in]{geometry}
\usepackage{amsmath,amssymb,mathtools}
\usepackage{graphicx}
\usepackage{booktabs}
\usepackage{hyperref}
\usepackage{xcolor}
\usepackage{enumitem}
\usepackage{algorithm}
\usepackage{algpseudocode}
\usepackage{caption}
\usepackage{url}
\usepackage{microtype}

\hypersetup{
  colorlinks=true,
  linkcolor=blue!60!black,
  citecolor=blue!60!black,
  urlcolor=blue!60!black
}

\title{The Brittleness of Plasticity:\\
A Synthesis for Neural Network Stability via Mask Evolution Operators}
\author{Paolo Pignatelli}
\date{September 2025}

% ---------- Macros ----------
\newcommand{\EWC}{\textsc{EWC}}
\newcommand{\MEO}{\textsc{MEO}}
\newcommand{\Id}{\mathrm{Id}}
\newcommand{\R}{\mathbb{R}}

\begin{document}
\maketitle

\begin{abstract}
We study the stability--plasticity trade-off in continual learning through \emph{Mask Evolution Operators} (\MEO s): task-agnostic transformations applied at activation level that exert a restoring force on drifting representations. We frame \MEO{} as a family of evolution operators over internal states and instantiate the \emph{Identity Evolution} case as a stringent stress test: representations are lightly regularized toward a reference manifold while the model continues to learn. 

We report smoke-test experiments on CIFAR-100 (10 tasks $\times$ 10 classes) using ResNet-50. On an Apple M1 (MPS backend), a full finetune baseline yields final average accuracy \textbf{6.28\%}; an \EWC{} baseline (best $\lambda$ from a fast grid, then confirmed at 20 epochs) yields \textbf{6.21\%}. These results are limited by short training schedules and CPU-class hardware; they serve as \emph{sanity checks} validating our pipeline and metrics prior to full GPU runs. Code and configs are available at \href{https://github.com/PaoloTCS/meo-continual-learning}{github.com/PaoloTCS/meo-continual-learning}. We position \MEO{} as a general operator framework; extending beyond Identity to controlled-evolution masks is left for future work.
\end{abstract}

\vspace{0.5em}
\noindent\textbf{Availability and Artifacts.} Code, configs, and scripts: \href{https://github.com/PaoloTCS/meo-continual-learning}{github.com/PaoloTCS/meo-continual-learning}. Figures are generated from the repository’s \texttt{results/} directory.

\section{Introduction}
Catastrophic forgetting remains a central obstacle in continual learning:\ as models adapt to new tasks, internal representations drift, eroding performance on earlier tasks. Classical mitigations constrain weights (\EWC{}, SI) or expand capacity (dynamic architectures), but each carries trade-offs in flexibility, compute, or memory. 

We propose a representation-centric view:\ \emph{Mask Evolution Operators} (\MEO s) act directly on activations, introducing a tunable, vector-field--like restoring force that discourages destructive drift while preserving useful plasticity. Conceptually, \MEO{} decouples \emph{where} learning occurs (the usual optimizer over weights) from \emph{how} state evolves in latent space (a mask-based operator). Practically, it is lightweight and architecture-agnostic.

This paper contributes: (i) a formalization of \MEO{} as an evolution operator over activations; (ii) an \emph{Identity Evolution} instantiation used as a \emph{stress test} to probe stability without freezing learning; (iii) a reproducible pipeline with drift metrics; (iv) smoke-test empirical results that validate the setup and motivate full GPU experiments.

\section{Related Work}
\textbf{Weight-space constraints.} \EWC{} penalizes deviations from Fisher-weighted optima~\cite{kirkpatrick2017overcoming}. Synaptic Intelligence and related approaches follow similar principles with different estimators. 
\textbf{Representation-space control.} Feature-level regularizers, contrastive consolidation, and orthogonality constraints modulate latent geometry directly. \MEO{} generalizes this idea by treating the controller as an explicit evolution operator. 
\textbf{Continual learning protocols.} Split CIFAR-100 and similar benchmarks remain standard for assessing drift and forgetting under multi-task streams.

\section{Mask Evolution Operators}
Let $h \in \R^d$ denote an activation vector at some layer. A \emph{Mask Evolution Operator} is a map
\begin{equation}
  \mathcal{M}_\theta: \R^d \to \R^d, \qquad \tilde{h} = \mathcal{M}_\theta(h),
\end{equation}
parameterized by $\theta$ (which may be learned or fixed), applied during the forward pass. We define a \emph{mask loss} that penalizes deviations from a reference state $h^\star$:
\begin{equation}
  \mathcal{L}_{\mathrm{mask}}(h; h^\star) \;=\; \alpha\,\| \mathcal{M}_\theta(h) - h^\star\|_2^2,
\end{equation}
with stiffness $\alpha \!>\! 0$ governing the restoring magnitude. The total training loss is
\begin{equation}
  \mathcal{L} \;=\; \mathcal{L}_{\mathrm{task}} + \mathcal{L}_{\mathrm{mask}}.
\end{equation}

\paragraph{Identity Evolution as stress test.}
In this paper we instantiate $\mathcal{M}_\theta$ as the identity $\Id$, anchoring representations to a fixed $h^\star$ captured early in training (Task~1). This deliberately favors stability and tests whether a light restoring force can deter destructive drift while still allowing useful adaptation. In future work we will allow $\mathcal{M}_\theta$ to evolve (e.g., low-rank adaptations, geodesic updates), enabling controlled motion of $h^\star$ across tasks.

\paragraph{Drift metric.}
We measure representation drift at layer $\ell$ via 
\begin{equation}
  \Delta_\ell \;=\; \mathbb{E}_{x \sim \mathcal{D}}\Big[\;\|h_\ell(x;t) - h_\ell^\star(x)\|_2\;\Big],
\end{equation}
comparing current activations to the reference. We report per-epoch and per-task aggregates.

\section{Baselines and Training}
\textbf{Finetune.} Standard ERM without rehearsal or consolidation.\\
\textbf{\EWC{}.} Online \EWC{} with Fisher estimated by gradient-squared averages over minibatches; we sweep $\lambda$ on a log scale and confirm the best at longer training.

\paragraph{Architecture and data.}
ResNet-50 on CIFAR-100 split into 10 tasks of 10 classes each; shared head with 100-way classification. Input preprocessing follows standard CIFAR normalization.

\paragraph{Hardware and schedule (smoke tests).}
All runs were executed on an Apple M1 (macOS, PyTorch MPS backend). For feasibility we used short schedules: 10 epochs per task for the grid, then 20 for the best \EWC{} $\lambda$. These are \emph{sanity checks} to validate the pipeline; full GPU results will follow.

\paragraph{Optimization and Fisher details.}
SGD with momentum $0.9$, weight decay $5\!\times\!10^{-4}$, cosine LR, batch size 128 (reduced if memory-bound). For \EWC{}, Fisher accumulation uses batch-averaged $g^2$ with online decay $\gamma=0.9$; we normalized by sample count to stabilize scale.

\section{Results}
Table~\ref{tab:main} summarizes final average accuracies (mean over the 10 tasks after the full stream). Numbers reflect the smoke-test setting on M1/MPS and should be interpreted as preliminary.

\begin{table}[h]
\centering
\begin{tabular}{lcc}
\toprule
Method & Final Avg. Acc. (\%) & Notes \\
\midrule
Finetune & \textbf{6.28} & 10 tasks $\times$ 10 epochs \\
\EWC{} (best $\lambda{=}30$) & \textbf{6.21} & grid@10 epochs, confirm@20 \\
\bottomrule
\end{tabular}
\caption{Smoke-test results on CIFAR-100 (ResNet-50) using Apple M1 (MPS). Values are expectedly low under short schedules and CPU-class hardware; they validate the pipeline prior to full GPU runs.}
\label{tab:main}
\end{table}

\paragraph{Drift and stiffness.}
Figure~\ref{fig:drift} shows per-epoch drift under a short schedule; Figure~\ref{fig:alpha} sweeps the stiffness $\alpha$ for \MEO{}. Both illustrate expected trends but should be considered qualitative under these constraints.

\section{Discussion}
\textbf{Identity as stress test.} A key critique of operator rigidity is addressed by positioning the Identity Evolution as a \emph{stress test}:\ it intentionally biases toward stability to probe whether even a simple operator can resist destructive drift without prohibiting adaptation. Our smoke tests indicate the pipeline works as intended; definitive conclusions require full GPU runs and non-identity operators.

\textbf{Beyond Identity.} The \MEO{} formalism naturally admits controlled-evolution operators (e.g., low-rank masks, proximal/geodesic updates) that move $h^\star$ along a constrained manifold. We hypothesize these will better preserve \emph{beneficial} plasticity while suppressing destructive drift.

\textbf{Limitations.} Short schedules, MPS backend, and absence of rehearsal constrain absolute performance. Our focus here is methodological clarity and reproducibility; scaling up is orthogonal and planned.

\section{Reproducibility}
Repository:\ \href{https://github.com/PaoloTCS/meo-continual-learning}{github.com/PaoloTCS/meo-continual-learning}. We provide configs, scripts for \EWC{} grid sweeping and paper auto-update, and plotting utilities for drift and $\alpha$-sweeps. The exact runs reported here used seed~42; see \texttt{results/} JSON logs.

\section*{Acknowledgements}
This research was conducted in a hybrid human--AI collaborative workflow. The author served as scientific lead, with \textbf{ChatGPT (GPT-5)} acting as the primary technical assistant and editor. Additional contributions came from \textbf{Claude}, \textbf{Grok}, and \textbf{Gemini}, which provided peer-review style feedback and alternative perspectives. This iterative, dialogic process reflects our broader R\&D philosophy:\ scientific progress as a conversational, multi-agent activity in which human insight and machine reasoning co-evolve.

% ---------- Algorithms / Pseudocode ----------
\begin{figure}[t]
\centering
\includegraphics[width=0.9\linewidth]{drift_placeholder.png}
\caption{\textbf{Drift per epoch (smoke test).} CIFAR-100, ResNet-50, Apple M1 (MPS), seed=42, short schedule. Higher-than-ideal magnitudes reflect limited training; the trend illustrates how the restoring term affects representation movement. Replace this placeholder with the generated figure when available.}
\label{fig:drift}
\end{figure}

\begin{figure}[t]
\centering
\includegraphics[width=0.9\linewidth]{alpha_sweep_placeholder.png}
\caption{\textbf{Stiffness sweep $\alpha$ (smoke test).} Validation of the \MEO{} stiffness parameter under short schedules. Mark the operating point used in the main experiment. Replace this placeholder with the generated sweep figure when available.}
\label{fig:alpha}
\end{figure}

% ---------- Optional pseudocode example ----------
\begin{algorithm}[t]
\caption{Training with \MEO{} (Identity Evolution stress test)}
\begin{algorithmic}[1]
\Require tasks $\{\mathcal{T}_k\}_{k=1}^K$, model $f_\phi$, reference activations $h^\star$ (from Task~1), stiffness $\alpha$
\For{$k=1$ to $K$}
  \For{minibatch $(x,y)\sim \mathcal{T}_k$}
    \State $h \gets \mathrm{forward\_activations}(f_\phi, x)$
    \State $\tilde{h} \gets h$ \Comment{Identity evolution}
    \State $\mathcal{L}_{\mathrm{task}} \gets \mathrm{CE}(f_\phi(x), y)$
    \State $\mathcal{L}_{\mathrm{mask}} \gets \alpha \|\tilde{h}-h^\star\|_2^2$
    \State $\mathcal{L} \gets \mathcal{L}_{\mathrm{task}} + \mathcal{L}_{\mathrm{mask}}$
    \State Update $\phi$ by SGD
  \EndFor
  \State Optionally update statistics (drift, Fisher, etc.)
\EndFor
\end{algorithmic}
\end{algorithm}

% ---------- References ----------
\begin{thebibliography}{9}

\bibitem{kirkpatrick2017overcoming}
Kirkpatrick, J., et al. Overcoming catastrophic forgetting in neural networks. \emph{PNAS}, 2017.

% Add your existing references here (or switch to your .bib file)
\end{thebibliography}

\end{document}